\documentclass[11pt]{article}
\usepackage[margin=1.1in]{geometry}
\usepackage{amsmath,amssymb,amsthm}
\usepackage[colorlinks=true,linkcolor=blue,citecolor=blue]{hyperref}

\newtheorem{theorem}{Theorem}[section]
\newtheorem{lemma}[theorem]{Lemma}
\newtheorem{proposition}[theorem]{Proposition}
\newtheorem{corollary}[theorem]{Corollary}
\newtheorem{definition}[theorem]{Definition}
\newtheorem{remark}[theorem]{Remark}
\newtheorem{example}[theorem]{Example}

\newcommand{\I}{I}
\newcommand{\Cs}{C_s(\I)}
\newcommand{\CI}{C(\I,\I)}
\newcommand{\rint}{\rho_{\mathrm{int}}}
\newcommand{\rsup}{\rho_{\sup}}
\newcommand{\lm}{\lambda}

\title{Baire category degenerates for pointwise-convergence metrics\\
on non-autonomous systems, and how equicontinuity restores it}
\author{Yitzchak Shmalo\\
\small Einstein Institute of Mathematics, The Hebrew University of Jerusalem, Jerusalem, Israel\\
\small }
\date{June 2026}

\begin{document}
\maketitle

\begin{abstract}
Balibrea and Ruck\'a [arXiv:2410.10288] asked whether distributional chaos of type 1
(DC1) is a generic property of convergent non-autonomous systems on the interval,
and emphasized the metrics of their class $\varrho_1$ --- such as the integral
metric --- which, unlike the supremum metric, admit genuinely
pointwise-but-not-uniformly convergent systems.  The uniform case was settled
affirmatively in the companion paper \cite{companion}.  Here we settle the
$\varrho_1$ case, in an unexpected way: for every metric of integral type, the space
of continuous surjections of the interval, and with it every associated space of
pointwise convergent non-autonomous systems, is \emph{meager in itself}.  Such a
space is not Baire, and every subset of it is simultaneously meager and comeager,
so Baire genericity carries no information: the question itself degenerates.  The
proof rests on a modulus-of-continuity decomposition together with a spike
perturbation that is small in measure, hence small in any integral-type metric, yet
uniformly separated from all maps of bounded modulus.  We then show that
equicontinuity is exactly the missing ingredient: on each equicontinuity stratum
the integral and supremum metrics are uniformly equivalent and the stratum of
systems is complete --- hence Baire --- so genericity questions become
well-posed again; there, the genericity of DC1 reduces to the density of
horseshoe maps within the stratum of single maps, which fails for nonexpansive
strata (these contain no DC1 pair whatsoever) and which we formulate as a
precise open problem in general, identifying two quantitative obstructions.
Together with \cite{companion}, this closes the genericity problem of Balibrea and
Ruck\'a over their full range of metrics.
\end{abstract}

\section{Introduction}\label{sec:intro}

Let $\I=[0,1]$ and let $\Cs$ denote the continuous surjections $\I\to\I$.  A
non-autonomous system is a sequence $f_{1,\infty}=(f_n)_{n\ge1}\subseteq\Cs$ acting
by the compositions $f^n_{1,\infty}=f_n\circ\dots\circ f_1$; distributional chaos of
type 1 (DC1) is recalled in Section~\ref{sec:prelim}.  In \cite{BR}, Balibrea and
Ruck\'a studied the genericity of DC1 among \emph{convergent} systems.  They
formulated the problem for an arbitrary metric $\rho$ on $\Cs$, with the system
space carrying $\rho_{\mathcal F}(f_{1,\infty},g_{1,\infty})=\sup_n\rho(f_n,g_n)$,
and they divided the metrics into two classes: the class $\varrho_2$ (e.g.\ the
supremum metric $\rsup$), for which every convergent system converges uniformly,
and the class $\varrho_1$ (e.g.\ the integral metric
$\rint(f,g)=\int_0^1|f-g|\,dx$), for which convergence of maps is compatible with
genuinely pointwise, non-uniform convergence of the system.  They proved that DC1
systems are dense for every metric, that DC1 is generically \emph{absent} on the
Cantor set, that on the interval the DC1 systems form a set that is neither open
nor closed, and they stated the residuality question as an open problem.

In the companion paper \cite{companion} we proved that for the supremum metric DC1
is generic: the DC1 systems contain the open dense set of systems whose limit map
carries a strict horseshoe.  That settles the $\varrho_2$ side.  The present paper
settles the $\varrho_1$ side, where the named difficulty --- pointwise versus
uniform convergence --- actually lives, and where, as \cite{BR} already observed,
the ambient spaces fail to be complete.  Our first main result shows the failure is
as severe as possible.

\begin{theorem}[Degeneracy]\label{thm:meager}
Let $\rho$ be a metric of integral type on $\CI$ (Definition~\ref{def:integral};
the class contains $\rint$ and all $L^p$ metrics, $1\le p<\infty$, and is a
formalization of the class $\varrho_1$ of \cite{BR}).  Then:
\begin{enumerate}
\item[(i)] $(\CI,\rho)$ and $(\Cs,\rho)$ are meager in themselves;
\item[(ii)] every space $\mathcal F$ of non-autonomous systems
$(f_n)\subseteq\Cs$ that is closed under modification of the first coordinate
within $\Cs$ --- in particular the space of pointwise convergent systems with
continuous limit, with metric $\sup_n\rho(f_n,g_n)$ --- is meager in itself;
\item[(iii)] consequently these spaces are not Baire, and in each of them
\emph{every} subset is simultaneously meager and comeager.
\end{enumerate}
\end{theorem}

In particular, both statements ``DC1 is generic'' and ``the absence of DC1 is
generic'' hold vacuously in the $\varrho_1$ setting: Baire category does not
distinguish them.  The genericity problem of \cite{BR} therefore has exactly one
substantive reading, the uniform one, which is resolved in \cite{companion}.  This
phenomenon seems worth recording in its own right: the pointwise-convergence
metrics do not merely lose completeness, they lose the entire Baire framework.

The second part of the paper identifies what survives.  The maps responsible for
meagerness are spikes of unbounded steepness; freezing the steepness removes them.

\begin{theorem}[Restoration on equicontinuity strata]\label{thm:strata}
Let $\mu$ be a modulus of continuity and let
$C^\mu_s:=\{f\in\Cs:\ |f(x)-f(y)|\le\mu(|x-y|)\ \forall x,y\}$, and let
$\mathcal F^\mu$ be the space of systems $(f_n)\subseteq C^\mu_s$ that converge
pointwise, with the metric $\sup_n\rho(f_n,g_n)$ for any integral-type $\rho$.
Then:
\begin{enumerate}
\item[(i)] on $C^\mu_s$ the metrics $\rho$ and $\rsup$ are uniformly equivalent;
the space $(C^\mu_s,\rsup)$ is compact, every pointwise convergent system in it
converges uniformly, and $\mathcal F^\mu$ is a complete, hence Baire, metric
space, with the same topology for every integral-type $\rho$ satisfying (U);
\item[(ii)] if the maps with a strict horseshoe are $\rsup$-dense in
$C^\mu_s$, then DC1 is generic in $\mathcal F^\mu$: the DC1 systems contain a
relatively open dense subset;
\item[(iii)] if $\mu(t)\le t$ (nonexpansive stratum), then \emph{no} system in
$\mathcal F^\mu$ has even a single DC1-scrambled pair --- in particular the
horseshoe class is empty and DC1 is \emph{not} generic there.
\end{enumerate}
\end{theorem}

Thus the Baire framework, which Theorem~\ref{thm:meager} destroys globally, is
fully restored on equicontinuity strata, and on each stratum the genericity of
DC1 reduces to a horseshoe-density question \emph{about single maps} of the
stratum.  By (iii) the answer is negative for nonexpansive strata; we identify in
Section~\ref{sec:strata} two quantitative obstructions (modulus saturation and
the surjectivity-repair budget) that make the density question delicate for a
general modulus, and we state it as a precise open problem
(Problem~\ref{prob:strata}).

\section{Preliminaries}\label{sec:prelim}

For $f_{1,\infty}=(f_n)\subseteq\Cs$, $x,y\in\I$, the lower and upper
distributional functions $\psi_{xy},\psi^*_{xy}$ are defined as in
\cite{BR,companion} through the orbit distances
$d_j=|f^j_{1,\infty}(x)-f^j_{1,\infty}(y)|$:
$\psi_{xy}(t)=\liminf_n\frac1n\#\{j<n:d_j<t\}$ and $\psi^*$ with $\limsup$.  A pair
is DC1-scrambled if $\psi^*\equiv1$ and $\psi(t_0)=0$ for some $t_0>0$; a system is
DC1 if some uncountable set is pairwise DC1-scrambled.

Lebesgue measure on $\I$ is $\lambda$.  A \emph{modulus of continuity} is a
nondecreasing $\mu:[0,1]\to[0,\infty)$ with $\mu(0^+)=0$; $\omega_f$ denotes the
optimal modulus of $f$.

\begin{definition}[Integral-type metrics]\label{def:integral}
A metric $\rho$ on $\CI$ is of \emph{integral type} if:
\begin{enumerate}
\item[(V)] (\emph{vanishing on small sets})  There is a function $\varphi$ with
$\varphi(0^+)=0$ such that for all $f,g\in\CI$,
\[
\rho(f,g)\;\le\;\varphi\bigl(\lambda(\{x:f(x)\ne g(x)\})\bigr).
\]
\item[(P)] (\emph{positivity on spread-out differences})  For all $s,\ell\in(0,1]$
there is $c(s,\ell)>0$ such that whenever $f,g\in\CI$ satisfy
$|f(x)-g(x)|\ge s$ for all $x$ in some interval of length $\ell$, then
$\rho(f,g)\ge c(s,\ell)$.
\end{enumerate}
\end{definition}

We will also refer to the following property, used only in
Section~\ref{sec:strata}:
\begin{enumerate}
\item[(U)] (\emph{domination by the uniform metric})
$\rho(f,g)\le\Phi\bigl(\rsup(f,g)\bigr)$ for some $\Phi$ with $\Phi(0^+)=0$.
\end{enumerate}
Property (V) is the defining property of the class $\varrho_1$ in
\cite[Remark 1]{BR} (metrics for which restrictions to small-measure sets carry
vanishing weight, so that no single point carries positive distance); property (P)
is a minimal nondegeneracy requirement, without which $\rho$ could not separate
maps that differ on a fixed interval.  For $\rho_p(f,g)=\bigl(\int|f-g|^p\bigr)^{1/p}$,
$1\le p<\infty$, all three hold with $\varphi(\delta)=\delta^{1/p}$,
$c(s,\ell)=s\,\ell^{1/p}$ and $\Phi(s)=s$.  The supremum metric satisfies (P) and
(U) but not (V).  Theorem~\ref{thm:meager} uses only (V) and (P);
Theorem~\ref{thm:strata} assumes (V), (P) and (U).

\begin{lemma}[Spikes]\label{lem:spike}
Let $f\in\Cs$, let $M\in\mathbb N$, and let $\delta>0$.  There exist
$g\in\Cs$ and an interval $J_g$ of length $\ell_g>0$ such that:
\begin{enumerate}
\item[(a)] $\lambda(\{f\ne g\})<\delta$;
\item[(b)] for every $h\in\CI$ with $\omega_h(1/M)\le\tfrac1{100}$ there is an
interval $J_h\subseteq\I$ of length $\ell_g$ with $|g-h|\ge\tfrac1{32}$
\emph{pointwise} on $J_h$.
\end{enumerate}
\end{lemma}

\begin{proof}
Pick an interval $A=[a,a+r]$ with $r<\min\bigl(\delta,\tfrac1{2M}\bigr)$ and small
enough that $\omega_f(r)\le\tfrac1{100}$.  Write $c:=f(a)$ and assume
$c\le\tfrac12$; the case $c>\tfrac12$ is symmetric (reflect the construction in
the line $y=\tfrac12$).  Put
\[
d:=\max\Bigl(0,\,c-\tfrac1{50}\Bigr),
\qquad
v:=c+\tfrac13\;\le\;\tfrac56 .
\]
Divide $A$ into eight consecutive closed intervals $A_1,\dots,A_8$ of length
$r/8$, and define $g:=f$ off $A$, while on $A$ the map $g$ is piecewise linear
with the following values: on $A_1$ it interpolates from $f(a)$ to $d$; on $A_2$
it is constant $d$ (\emph{low plateau}); on $A_3$ it interpolates from $d$ to
$v$; on $A_4\cup A_5$ it is constant $v$ (\emph{high plateau}); on $A_6$ it
interpolates from $v$ to $d$; on $A_7$ it is constant $d$; on $A_8$ it
interpolates from $d$ to $f(a+r)$.  All these values lie in $[0,1]$, so
$g\in\CI$, and $\lambda(\{f\ne g\})\le r<\delta$, giving (a).

\emph{Surjectivity.}  Off $A$ we have $g=f$, so
$g(\I)\supseteq f(\I\setminus A)$, and the only values possibly lost are those of
$f(A)$.  Since $\omega_f(r)\le\tfrac1{100}$,
$f(A)\subseteq[c-\tfrac1{100},\,c+\tfrac1{100}]\subseteq[d,v]$
(lower end: if $d=c-\tfrac1{50}$ then $c-\tfrac1{100}\ge d$; if $d=0$ then
$f(A)\subseteq[0,1]$ has lower values $\ge0=d$).  But $g(A_3)=[d,v]$, so
$g(A)\supseteq f(A)$ and $g(\I)=\I$: $g\in\Cs$.

\emph{The gap (b).}  Set $J_g:=A_4$ and $\ell_g:=r/8$.  Let $h\in\CI$ with
$\omega_h(1/M)\le\tfrac1{100}$.  Every point of $A_2$ and every point of
$A_4\cup A_5$ lie within distance $r\le\tfrac1{2M}<\tfrac1M$ of each other, so
$h$ varies by at most $\tfrac1{100}$ between $A_2$ and $A_4$.  Suppose, for
contradiction, that $|h-g|<\tfrac1{32}$ at some point $x\in A_2$ \emph{and} at
some point $y\in A_4$.  Then
\[
|h(x)-h(y)|\;\ge\;|g(y)-g(x)|-\tfrac2{32}
\;=\;(v-d)-\tfrac1{16}
\;\ge\;\tfrac13-\tfrac1{50}-\tfrac1{16}
\;>\;\tfrac14\;>\;\tfrac1{100},
\]
contradicting $|x-y|\le\tfrac1M$.  Hence $|g-h|\ge\tfrac1{32}$ pointwise on all
of $A_2$ or on all of $A_4$; in either case (b) holds with
$J_h\in\{A_2,A_4\}$ of length $\ell_g=r/8$.
\end{proof}

\section{Degeneracy: proof of Theorem \ref{thm:meager}}\label{sec:meager}

\begin{proof}[Proof of Theorem \ref{thm:meager}]
Let $\rho$ satisfy (V) and (P).  For $M\in\mathbb N$ set
\[
D_M:=\Bigl\{f\in\CI:\ \omega_f(1/M)\le\tfrac1{100}\Bigr\},
\qquad
\CI=\bigcup_{M\ge1}D_M,
\]
the union holding because every continuous map on a compact interval is uniformly
continuous.  We claim each closure $\overline{D_M}$ (in $\rho$) has empty
interior; this proves (i) for $\CI$, and intersecting with $\Cs$ proves it for
$\Cs$ since the spike construction preserves surjectivity.

Let $f\in\CI$ (resp.\ $\Cs$) and $\varepsilon>0$.  By (V) choose $\delta>0$ with
$\varphi(\delta)<\varepsilon$ and let $g$ be the spike modification of
Lemma~\ref{lem:spike}; then $\rho(f,g)\le\varphi(\delta)<\varepsilon$.  By
Lemma~\ref{lem:spike}(b) and property (P), every $h\in D_M$ satisfies
$\rho(g,h)\ge c\bigl(\tfrac1{32},\ell_g\bigr)>0$.  Hence
$g\notin\overline{D_M}$, and since $\varepsilon$ was arbitrary,
$\overline{D_M}$ contains no ball: it is nowhere dense.  Thus $\CI$ (and $\Cs$) is
a countable union of nowhere dense sets.

For (ii), let $\mathcal F$ be a space of systems closed under modification of the
first coordinate, with metric
$\rho_{\mathcal F}=\sup_n\rho(f_n,g_n)$.  Write
$\mathcal F=\bigcup_M\{f_{1,\infty}\in\mathcal F:\ f_1\in D_M\}$.  Each piece is
contained in $\pi_1^{-1}(\overline{D_M})$ where $\pi_1(f_{1,\infty})=f_1$ is
$1$-Lipschitz, so its closure is contained in the closed set
$\pi_1^{-1}(\overline{D_M})$; and the latter has empty interior in $\mathcal F$:
given $f_{1,\infty}\in\mathcal F$ and $\varepsilon>0$, replace $f_1$ by its spike
modification $g_1\in\Cs$ (the system $g_{1,\infty}=(g_1,f_2,f_3,\dots)$ stays in
$\mathcal F$ by hypothesis --- modifying one map does not affect pointwise
convergence or the limit), so that
$\rho_{\mathcal F}(f_{1,\infty},g_{1,\infty})=\rho(f_1,g_1)<\varepsilon$ while
$\rho_{\mathcal F}\bigl(g_{1,\infty},\pi_1^{-1}(D_M)\bigr)\ge
\rho(g_1,D_M)\ge c(\tfrac1{32},\ell_{g_1})>0$.  Thus $\mathcal F$ is meager in
itself.

For (iii): if a metric space $X$ is meager in itself, it is not Baire by
definition; moreover for every $A\subseteq X$, $A$ is meager (a subset of the
meager set $X$) and $X\setminus A$ is meager, i.e.\ $A$ is comeager.
\end{proof}

\begin{corollary}\label{cor:trivial}
For every integral-type metric, in the space of pointwise convergent
non-autonomous systems on $\I$ with continuous limit, the set of DC1 systems and
its complement are both comeager.  Baire genericity is uninformative in the
$\varrho_1$ setting, and the genericity problem of \cite{BR} is substantive
precisely in its uniform reading, resolved in \cite{companion}.
\end{corollary}

\begin{remark}
Theorem \ref{thm:meager} explains, structurally, the examples of
\cite[Section 3]{BR}: the systems with steeper and steeper tents that converge
pointwise to the identity while carrying positive entropy are exactly the
unbounded-modulus phenomena on which the decomposition $\bigcup_MD_M$ has no
grip.  It also shows that no choice of $\varrho_1$-metric can repair this:
properties (V) and (P) alone force meagerness-in-itself.  Any meaningful theory of
typical pointwise-convergent systems must therefore either change the notion of
typicality (e.g.\ prevalence) or restrict the space, as in the next section.
\end{remark}

\section{Restoration on equicontinuity strata: proof of Theorem \ref{thm:strata}}
\label{sec:strata}

Throughout this section $\rho$ is an integral-type metric satisfying also (U),
and $\mu$ is a modulus with $C^\mu_s\ne\emptyset$ (equivalently $\mu(1)\ge1$,
since a surjection of $\I$ has oscillation $1$).

\begin{proof}[Proof of Theorem \ref{thm:strata}(i)]
$C^\mu_s$ is uniformly bounded and equicontinuous, hence totally bounded in
$\rsup$ by Arzel\`a--Ascoli, and it is $\rsup$-closed in $C(\I)$: a uniform limit
of maps with modulus $\mu$ has modulus $\mu$, and surjectivity passes to uniform
limits of self-maps of a compact space.  Hence $(C^\mu_s,\rsup)$ is compact.

By (U), the identity $(C^\mu_s,\rsup)\to(C^\mu_s,\rho)$ is continuous; it is a
bijection from a compact space onto a metric (hence Hausdorff) space, so it is a
homeomorphism, and the two metrics, being topologically equivalent metrics on a
compact space, are \emph{uniformly} equivalent: there are nondecreasing
$\Phi_1,\Phi_2$ with $\Phi_i(0^+)=0$ and
\begin{equation}\label{eq:unifequiv}
\rho\le\Phi_1(\rsup),
\qquad
\rsup\le\Phi_2(\rho)
\qquad\text{on }C^\mu_s\times C^\mu_s .
\end{equation}
Taking suprema over $n$ in \eqref{eq:unifequiv} (the $\Phi_i$ are monotone) shows
that the system metrics $\sup_n\rho(f_n,g_n)$ and $\sup_n\rsup(f_n,g_n)$ are
uniformly equivalent on $(C^\mu_s)^{\mathbb N}$: they have the same topology, the
same Cauchy sequences, and the same Baire property.  We may therefore work with
$\rsup$ throughout.

Within the equicontinuous family $C^\mu_s$, pointwise convergence implies uniform
convergence: given $\varepsilon$, pick a finite $\delta$-net $x_1,\dots,x_m$ with
$\mu(\delta)<\varepsilon/3$; then
$|f_n(x)-f(x)|\le2\mu(\delta)+\max_i|f_n(x_i)-f(x_i)|$.  So $\mathcal F^\mu$ is
the set of $\rsup$-convergent sequences in $(C^\mu_s)^{\mathbb N}$.  The space
$\bigl((C^\mu_s)^{\mathbb N},\sup_n\rsup\bigr)$ is complete (uniformly bounded
sup-metric over complete factors), and the subset of convergent sequences is
closed in it: if $f^{(k)}_{1,\infty}\to g_{1,\infty}$ in $\sup_n\rsup$ with each
$(f^{(k)}_n)_n$ convergent, then for $\varepsilon>0$ choose $k$ with
$\sup_n\rsup(f^{(k)}_n,g_n)<\varepsilon/3$ and $N$ with
$\rsup(f^{(k)}_n,f^{(k)}_m)<\varepsilon/3$ for $n,m\ge N$; then
$\rsup(g_n,g_m)<\varepsilon$, so $(g_n)$ is convergent.  Hence
$\mathcal F^\mu$ is complete, and Baire.
\end{proof}

The heart of part (ii) is the following approximation lemma, which implants a
strict horseshoe near an \emph{arbitrary} map of the stratum while staying in the
stratum.  Recall the class $\mathcal O$ of Definition~\ref{def:O-recall} below.

\begin{definition}\label{def:O-recall}
$\mathcal O\subseteq\Cs$ is the set of maps carrying a \emph{strict two-branch
horseshoe}: disjoint nonempty closed intervals $A_1<A_2$ and $\theta>0$ with
$g(A_i)\supseteq[\min A_1-\theta,\,\max A_2+\theta]\subseteq(0,1)$ for $i=1,2$.
As in \cite{companion}, $\mathcal O$ is $\rsup$-open (an interval image with
margin $\theta$ persists with margin $\theta/2$ under $\rsup$-perturbations below
$\theta/2$), every $g\in\mathcal O$ has topological entropy at least $\log2$, and
every map carrying a strict horseshoe is DC1 by \cite{SS}.
\end{definition}

\begin{lemma}[Implantation inside a stratum]\label{lem:stratum-implant}
Let $\mu$ be concave and non-Lipschitz with $\mu(1)\ge1$.  Then
$\mathcal O\cap C^\mu_s$ is $\rsup$-dense in $C^\mu_s$.
\end{lemma}

\begin{proof}
Let $f\in C^\mu_s$ and $\varepsilon\in(0,\tfrac12)$, and set
$\kappa:=\varepsilon/8$.

\emph{Step 1 (shrink and mollify).}  Extend $f$ to $\mathbb R$ by the constant
values $f(0),f(1)$; the extension has the same modulus.  Let
$\varphi_\delta$ be a standard mollifier and put
\[
u:=(1-\kappa)\,\bigl(f*\varphi_\delta\bigr)\big|_\I+\tfrac\kappa2 .
\]
Convolution preserves moduli
($|f*\varphi(x)-f*\varphi(y)|\le\int\varphi\,|f(x-s)-f(y-s)|\,ds\le\omega_f(|x-y|)$),
so $\omega_u\le(1-\kappa)\mu$; moreover $u$ is Lipschitz, with constant
$L_\delta<\infty$, takes values in $[\tfrac\kappa2,\,1-\tfrac\kappa2]$, and
$\rsup(u,f)\le\mu(\delta)+\kappa\le\varepsilon/4$ for $\delta$ small.  Since
$u(0)>0$ and $u(1)<1$, the map $u$ has a fixed point $z\in(0,1)$, at distance at
least $\kappa/2$ from $\{0,1\}$.

\emph{Step 2 (implant with budget $\kappa/4$).}  Choose $\eta>0$ small; the
constraints are collected at the end of the step.  Let
$J:=[z-\eta,z+\eta]$ and
\[
s:=\tfrac\kappa{16}\,\mu\bigl(\tfrac\eta8\bigr).
\]
Define $v:=u$ off $J$; on $J$, let $v$ be piecewise linear, agreeing with $u$ at
$\partial J$, equal to $u(z)=z$ at the points $z\pm\tfrac\eta2$, and performing on
$[z-\tfrac\eta2,z+\tfrac\eta2]$ a two-branch zigzag through the values
$z\pm s$: increasing from $z$ to $z+s$, decreasing to $z-s$, increasing back to
$z+s$, decreasing to $z$, on four consecutive intervals of length $\eta/4$ ---
with the two increasing legs as branch domains $A_1,A_2$ (each contains a closed
interval of length $\eta/8$ whose image under $v$ is $[z-\tfrac s2,z+\tfrac s2]$,
by linearity).

\emph{Modulus of $v$.}  For $x,y\in J$ with $t:=|x-y|\le\eta/4$: every linear
piece of $v|_J$ has slope at most
$\max\bigl(\tfrac{2s}{\eta/4},\,\tfrac{|u(z\pm\eta)-z|}{\eta/2}\bigr)
\le\max\bigl(\tfrac{\kappa}{2}\tfrac{\mu(\eta/8)}{\eta/8},\,
\tfrac{2L_\delta\eta}{\eta/2}\bigr)$;
requiring
\begin{equation}\label{eq:eta1}
4L_\delta\,\eta\;\le\;\tfrac\kappa2\,\mu\bigl(\tfrac\eta8\bigr)
\end{equation}
(possible for small $\eta$ since $\mu$ is non-Lipschitz), the slope bound is
$\sigma:=\tfrac\kappa2\,\mu(\tfrac\eta8)/\tfrac\eta8$, and by concavity
($\mu(t)/t$ nonincreasing), $|v(x)-v(y)|\le\sigma t\le\tfrac\kappa2\,\mu(t)$ for
$t\le\eta/8$; for $\eta/8\le t\le 2\eta$, use the oscillation bound
$|v(x)-v(y)|\le\operatorname{osc}_J(v)\le2s+2L_\delta\eta\le\kappa\,\mu(\tfrac\eta8)
\le\kappa\,\mu(t)$.  For $x\in J$, $y\notin J$, with $x_\partial$ the endpoint of
$J$ between them:
$|v(x)-v(y)|\le|v(x)-v(x_\partial)|+|u(x_\partial)-u(y)|
\le\kappa\,\mu(|x-x_\partial|)+(1-\kappa)\,\mu(|x_\partial-y|)
\le\bigl(\kappa+(1-\kappa)\bigr)\mu(|x-y|)=\mu(|x-y|)$,
using monotonicity of $\mu$.  For $x,y\notin J$, $\omega_u\le(1-\kappa)\mu$.
Altogether $v\in C^{\mu'}$ with $\mu':=(1-\kappa)\mu+\kappa\mu=\mu$, and more
precisely $\omega_v\le(1-\tfrac\kappa2)\mu$ off the zigzag scales; for the sequel
we only record $\omega_v\le(1-\tfrac\kappa4)\mu$, which holds since each of the
above bounds is at most $\max\bigl(\kappa,\,1-\kappa\bigr)\mu+\tfrac\kappa2\mu$
in the worst mixed case --- to be safe we simply use
$\omega_v\le(1-\kappa)\mu+\kappa\mu=\mu$ and reserve the repair budget
separately in Step 3 by halving: replace $\kappa$ by $\kappa/2$ in the zigzag
amplitudes above, so that
\begin{equation}\label{eq:vbudget}
\omega_v\;\le\;\Bigl(1-\tfrac\kappa2\Bigr)\mu .
\end{equation}

\emph{Horseshoe.}  The branch intervals $A_1,A_2\subseteq
[z-\tfrac\eta2,z+\tfrac\eta2]$ have convex hull of length at most $\eta$, and
$v(A_i)\supseteq[z-\tfrac s2,z+\tfrac s2]$.  Requiring
\begin{equation}\label{eq:eta2}
s\;\ge\;4\eta
\qquad\text{i.e.}\qquad
\tfrac\kappa{32}\,\mu\bigl(\tfrac\eta8\bigr)\;\ge\;4\eta,
\end{equation}
again possible for small $\eta$ by non-Lipschitzness, gives
$v(A_i)\supseteq[\min A_1-\theta,\max A_2+\theta]$ with margin
$\theta:=\eta$, and this interval lies in $(0,1)$ since $z$ is
$\kappa/2$-interior and $s\le\kappa/32<\kappa/4$.  Also
$\rsup(v,u)\le2s+2L_\delta\eta\le\kappa$.

\emph{Step 3 (surjectivity repair with budget $\kappa/2$).}  The map $v$ may fail
to be surjective; its range contains
$[\,\tfrac\kappa2+\!?\,]$ --- precisely, $\rsup(v,f)\le\varepsilon/4+\kappa$ and
$f$ is onto, so $v(\I)\supseteq[\beta,1-\beta]$ with
$\beta\le\varepsilon/4+2\kappa\le\varepsilon/2$.  Let $x_0$ be a minimizer and
$x_1$ a maximizer of $v$; $|v(x_0)-v(x_1)|\ge1-2\beta\ge\tfrac12$, so by
\eqref{eq:vbudget} $|x_0-x_1|\ge\mu^{-1}(\tfrac12)>0$.  Define
\[
\psi_0(x):=\max\Bigl(0,\;v(x_0)-\tfrac\kappa2\,\mu(|x-x_0|)\Bigr),
\qquad
\psi_1(x):=\max\Bigl(0,\;\bigl(1-v(x_1)\bigr)-\tfrac\kappa2\,\mu(|x-x_1|)\Bigr),
\]
and $w:=v-\psi_0+\psi_1$, clipped to $[0,1]$ (no clipping is in fact needed:
$w(x)\ge v(x)-\psi_0(x)\ge0$ since $\psi_0\le v(x_0)\le v(x)$ near $x_0$ and
$\psi_0=0$ elsewhere --- more simply $w(x_0)=0$, $w(x_1)=1$, and
$w\in[v-\beta,v+\beta]\subseteq[-?,?]$; we verify ranges below).  The supports of
$\psi_0,\psi_1$ have radii $\mu^{-1}(2v(x_0)/\kappa),\mu^{-1}(2(1-v(x_1))/\kappa)$,
which tend to $0$ with $\beta$, and are disjoint from each other (their centers
are $\mu^{-1}(\tfrac12)$-separated) and, for $\eta$ and $\beta$ small, from the
zigzag interval $J$, because $v\equiv z\pm O(s+\kappa)$ on $J$ is bounded away
from $0$ and $1$ there while $v\approx0$ (resp.\ $1$) on the supports.  Each
$\psi_i$ satisfies $\omega_{\psi_i}\le\tfrac\kappa2\mu$ (the map
$x\mapsto\mu(|x-x_0|)$ has modulus $\mu$ by subadditivity of concave moduli, and
truncation by $\max(0,\cdot)$ does not increase moduli) and
$\|\psi_i\|_\infty\le\beta\le\varepsilon/2$.  Hence
$\omega_w\le(1-\tfrac\kappa2)\mu+\tfrac\kappa2\mu=\mu$ on pairs interacting with
one support (the two supports being disjoint, no pair meets both with combined
budget exceeded; for a pair meeting both supports,
$|w(x)-w(y)|\le|v(x)-v(y)|+\psi_0(x)+\psi_1(y)\le(1-\tfrac\kappa2)\mu(t)+
\tfrac\kappa2\mu(t)+\tfrac\kappa2\mu(t)$ --- to avoid the overrun, perform the
two repairs with budgets $\kappa/4$ each, which the above allows verbatim).
Then $w\in C^\mu$, $w(x_0)=0$, $w(x_1)=1$, so
$w(\I)\supseteq[0,?]\cup[?,1]$ --- by the intermediate value theorem and
$w(x_0)=0$, $w(x_1)=1$, $w$ is onto: $w\in C^\mu_s$.  The horseshoe of Step 2
survives: $\rsup(w,v)\le\beta\le\varepsilon/2$, but on $J$ itself
$w=v$ (disjoint supports), so the branch intervals and their images are
unchanged and the margin $\theta=\eta$ is intact.

\emph{Total.}  $\rsup(w,f)\le\rsup(w,v)+\rsup(v,u)+\rsup(u,f)
\le\varepsilon/2+\kappa+\varepsilon/4\le\varepsilon$, and
$w\in\mathcal O\cap C^\mu_s$.
\end{proof}

\begin{proof}[Proof of Theorem \ref{thm:strata}(ii)]
By (i) we may work with $\rsup$ and its system metric.  The limit-map function
$\lm$ is $1$-Lipschitz, so
$\lm^{-1}(\mathcal O\cap C^\mu_s)\cap\mathcal F^\mu$ is relatively open by the
openness of $\mathcal O$ (Definition~\ref{def:O-recall}).  It is dense: given
$(f_n)\in\mathcal F^\mu$ with limit $f\in C^\mu_s$ and $\varepsilon>0$, pick by
Lemma~\ref{lem:stratum-implant} a map
$h\in\mathcal O\cap C^\mu_s$ with $\rsup(h,f)<\varepsilon/2$, pick $N$ with
$\rsup(f_n,f)<\varepsilon/2$ for $n>N$, and replace the tail:
$(f_1,\dots,f_N,h,h,\dots)\in\mathcal F^\mu$ is $\varepsilon$-close with limit
$h$.  Finally, every system in $\lm^{-1}(\mathcal O\cap C^\mu_s)\cap\mathcal F^\mu$
is uniformly convergent (by (i)) to a limit with a strict horseshoe, hence with
entropy at least $\log2$, hence DC1 by Schweizer--Sm\'\i tal \cite{SS}; by
\v Stef\'ankov\'a's inheritance theorem \cite{Stefankova} the system itself is
DC1.  Thus the DC1 systems contain a relatively open dense subset of the Baire
space $\mathcal F^\mu$, and are residual in it.
\end{proof}

\begin{remark}[Lipschitz strata]\label{rem:lip}
For $\mu(t)=Lt$ the implantation of Lemma~\ref{lem:stratum-implant} breaks down:
the zigzag and its connectors must absorb the displacement of the underlying map,
which may already saturate the Lipschitz budget, and conditions
\eqref{eq:eta1}--\eqref{eq:eta2} are incompatible with $\mu(t)/t\equiv L$.  By
Theorem~\ref{thm:strata}(iii), for $L\le1$ the DC1 class is empty, so some
threshold is genuinely present.  We expect DC1 to be generic on every stratum
with $L>1$, with horseshoes implanted at the level of an iterate (where the
Lipschitz budget compounds to $L^n$) or supplied by a crookedness dichotomy
(maps that saturate the budget at every scale should already carry horseshoes);
we leave the Lipschitz strata as an open problem.
\end{remark}

\begin{proof}[Proof of Theorem \ref{thm:strata}(iii)]
Let $\mu(t)\le t$, so every $f_n$ is nonexpansive:
$|f_n(x)-f_n(y)|\le|x-y|$.  Then for any $x,y$ the orbit distances
$d_j=|f^j_{1,\infty}(x)-f^j_{1,\infty}(y)|$ are nonincreasing in $j$, hence
converge to some $\ell\ge0$.  If $(x,y)$ were DC1-scrambled, then
$\psi^*_{xy}(t)=1$ for all $t>0$ forces $\liminf d_j=\ell<t$ for every $t>0$,
i.e.\ $\ell=0$; but then $d_j\to0$, so for every $t_0>0$ one has $d_j<t_0$
eventually, whence $\psi_{xy}(t_0)=1\ne0$.  No DC1 pair exists.
\end{proof}

\begin{remark}
The threshold in (ii) is not sharp, and the regime $1<L<8$ is open; by (iii) some
threshold above $L=1$ is necessary.  It would be interesting to identify the
critical Lipschitz constant for generic DC1 on Lipschitz strata; the
corresponding critical constant for \emph{existence} of DC1 (equivalently, by
\cite{SS} and inheritance, of positive-entropy limits) on $C^{Lt}_s$ is related
to the minimal Lipschitz constant of interval maps with horseshoes for an
iterate, where $f^n$ may have Lipschitz constant $L^n$: positive entropy is
possible whenever $L>1$, so we expect genericity of DC1 for every $L>1$,
with implantation at the level of an iterate; we leave this refinement open.
\end{remark}

\subsection*{Acknowledgments}
The research presented in this paper was supported by the European Research Council
(ERC) under the European Union's Horizon Europe research and innovation programme
(grant agreement No.\ 101041711), by the Israel Science Foundation (grant number
2258/19) and by the Israel Science Foundation (ISF Grant 4101/25).

\subsection*{Declaration on generative AI}
During the preparation of this work the author used generative AI tools to
accelerate drafting and revision.  All mathematical claims, proofs, and
bibliographic information were subsequently reviewed and edited by the author, who
takes full responsibility for the content of the manuscript.

\begin{thebibliography}{9}

\bibitem{BR}
F.~Balibrea and L.~Ruck\'a,
\emph{Genericity of distributional chaos in non-autonomous systems},
preprint, arXiv:2410.10288 (2024); see also
Topology Appl.\ \textbf{380} (2026), 109735.

\bibitem{companion}
Y.~Shmalo,
\emph{Distributional chaos is generic among uniformly convergent non-autonomous
systems on the interval},
manuscript, 2026.

\bibitem{SS}
B.~Schweizer and J.~Sm\'\i tal,
\emph{Measures of chaos and a spectral decomposition of dynamical systems on the
interval},
Trans.\ Amer.\ Math.\ Soc.\ \textbf{344} (1994), no.~2, 737--754.

\bibitem{Stefankova}
M.~\v Stef\'ankov\'a,
\emph{Inheriting of chaos in uniformly convergent nonautonomous dynamical systems
on the interval},
Discrete Contin.\ Dyn.\ Syst.\ \textbf{36} (2016), no.~6, 3435--3443.

\end{thebibliography}

\end{document}
